%%%%%%%%%%%%%%%%%% NOTE %%%%%%%%%%%%%%%%%%
% - This is a beta version of the University of Manchester LaTeX thesis template that adds the \DocumentMetadata command. This template produces PDF files which include a range of features to improve accessibility. If interested you can read about it here: https://www.latex-project.org/publications/indexbytopic/pdf/
% - The template fundamentally works, but produces a few warnings, and you may find that a number of third party packages don't work. (For example we usually use \tabularx for making fixed width tables, but it isn't compatible with DocumentMetadata with the current version in Overleaf.)
% - If you have a relatively 'vanilla' thesis, with just text, figures, and simple tables, everything will likely work with no issues. If this is the case, I would suggest you use this template. You get a PDF that works better with screen readers and other accessibility tools, for no extra effort. If you have issues using this template version though, then just use the 'regular' one at https://www.overleaf.com/latex/templates/university-of-manchester-phd-thesis-and-continuation-reports/hgngpshkptbd. There are some parts in the example which are just commented out because they don't work yet.
% - The intention is that this version of the template for producing more accessible PDF files will become the default template in the future, but it likely needs a few more years to become stable. Accessibility features are still being actively developed in LaTeX, and as Overleaf uses TexLive and packages LaTeX once per year, there can be quite a delay before new features become available. You can use it now if you want to though.
% - A number of items are flagged as "% TAGGING:" where known improvements/fixes are needed.



%%%%%%%%%%%%%%%%%% USAGE INSTRUCTIONS %%%%%%%%%%%%%%%%%%
% - Compile using LuaLaTeX and biber. Do not use the older LaTex/PDFLaTeX or BibTeX (the fonts won't work correctly).
% - Expect to get a number of warnings even when working correctly. These can be ignored.
% - Font, font size, and the report 'year' must all be specified in \documentclass or the template won't work correctly. (There's no error checking/default cases!)
% - Options for fonts are: calibri, times, palatino, garamond, arial, tahoma, verdana, trebuchet. Note however that not all these are installed on Overleaf, you need to install to your project if use (e.g. calibri)
% - Options for font size are: 10pt, 11pt, 12pt or 14pt. 12pt is the reccomended value in the thesis guidelines
% - Options for 'year' are: first, second, third, thesis.
% - Options for 'margin' are: online and print. If printing, it is reccomended there is a larger margin on the left side to allow for page turns, but this is no longer mandatory. The 'print' option sets this margin to 4 cm and all other margins to 2 cm (slightly larger than the minimum in the regulations of 1.5 cm). online sets all of the margins to 2 cm. 
% - As many further packages as wanted can be loaded. Below are just an example set. Note that template itself loads a number of packages, including hyperref and csquote.
% - Can use \uomaddblankpage if you want a blank page with blank page written on it
% - References are handed using biblatex.
% - There is more documentation in the header of uom_these_casson.cls file if you need more help
% - Link to the presentation of theses policy: http://www.regulations.manchester.ac.uk/pgr-presentation-theses/
% - This template has been checked to be compliant with the requirements. Nevertheless, responsibility for ensuring compliance with the University of Manchester Presentation of Theses Policy remains with the candidate.




%%%%%%%%%%%%%%%%%% DOCUMENT SETUP %%%%%%%%%%%%%%%%%%
\DocumentMetadata{lang=en-GB,pdfversion=2.0,pdfstandard=A-4F,pdfstandard=UA-2,testphase={phase-III,title,math,table,firstaid}} % edit language if needed, but otherwise leave unchanged
% TAGGING: will become this when Overleaf catches up with current LaTeX versions
%\DocumentMetadata{lang=en-GB,pdfversion=2.0,pdfstandard={UA-2,A-4F},tagging=on} 

\documentclass[times,12pt,thesis,online]{uom_thesis_casson} % See above for font options. Size can be: 10pt, 11pt, 12pt or 14pt. Year can be: first, second, third, or thesis. Margins can be set for: online or print. 


%%%%%%%%%%%%%%%%%% PACKAGES AND COMMANDS %%%%%%%%%%%%%%%%%%

% Packages - some useful examples
\usepackage{graphicx}              % for adding graphics files
  \graphicspath{ {./images/} }
\usepackage{amsmath}               % assumes amsmath package installed
  \allowdisplaybreaks[1]           % allow eqnarrays to break across pages
\usepackage{amssymb}               % assumes amsmath package installed 
\usepackage{url}                   % format hyperlinks correctly
\usepackage{rotating}              % allow portrait figures and tables
\usepackage{multirow}              % allows merging of rows in tables
\usepackage{lscape}                % allows pages to be typeset in landscape mode
% TAGGING: Re-add when supported for tagging
%\usepackage{tabularx}              % allows fixed width tables
\usepackage{verbatim}              % enhanced version of built-in verbatim environment
\usepackage{footnote}              % allows more control over footnote environments
\usepackage{float}                 % allows H option on floats to force here placement
\usepackage{booktabs}              % improve table line spacing
\usepackage{lipsum}                % for adding dummy text here
\usepackage[base]{babel}           % required for lisum package
\usepackage{subcaption}            % for multiple sub-figures in a single float
\usepackage{siunitx}               % add SI units
% Add your packages here


% Custom commands
\newcommand{\degree}{\ensuremath{^\circ}}
\newcommand{\sus}[1]{$^{\mbox{\scriptsize #1}}$} % superscript in text (e.g. 1st can be 1\sus{st})
\newcommand{\sub}[1]{$_{\mbox{\scriptsize #1}}$} % subscript in text
\newcommand{\otoprule}{\midrule[\heavyrulewidth]}
% TAGGING: Re-add when tabularx supported for tagging
%\newcolumntype{Z}{>{\centering\arraybackslash}X}  % tabularx centered columns 


% Add your custom commands here



%%%%%%%%%%%%%%%%%% REFERENCES SETUP %%%%%%%%%%%%%%%%%%
% Setup your references here. Change the reference style here if wanted
\usepackage[style=ieee,backend=biber,backref=true,hyperref=auto,maxbibnames=3,minbibnames=1]{biblatex}
% Note backref=true adds a page number (and hyperlink) to each reference so you can easily go back from the references to the main document. You may prefer backref=false if you need to stick strictly to a given reference style


% Fixes which can't be applied in the .cls file. Don't change
\DefineBibliographyStrings{english}{backrefpage = {cited on p\adddot},  backrefpages = {cited on pp\adddot}}
  % TAGGING: Check for different font sizes, and move away from extsizes
  \renewcommand*{\bibfont}{\normalsize}


% Add more .bib files here if wanted
\addbibresource{references.bib}



%%%%%%%%%%%%%%%%%% AUTOMATIC WORD COUNT SETUP %%%%%%%%%%%%%%%%%%
% Automatically counts the words present and makes a \mywordcount variable. This is used later to automatically add the word count (which can be overridden if wanted)
% See https://www.overleaf.com/learn/how-to/Is_there_a_way_to_run_a_word_count_that_doesn%27t_include_LaTeX_commands%3F for info on what words are counted and how to control this. Any changes to what's counted need to be made in uom_thesis_casson.cls, in the \newcommand{\quickwordcount} command. The default doesn't count words in headings, captions, or references and so you may get slightly different numbers if use a different word count tool
% This can be a bit fragile. If it doesn't work, there's an option later on to just type in a numer
\quickwordcount{\currfilebase} % run word count. This just counts the words. Is displayed with the \wordcount command later



%%%%%%%%%%%%%%%%%% START DOCUMENT %%%%%%%%%%%%%%%%%%
\begin{document}



%%%%%%%%%%%%%%%%%% LIST OF THESIS REVISIONS %%%%%%%%%%%%%%%%%%
%TC:ignore % turns off the word count for this section
%\begin{uomlotr} % only for resubmitted theses, if wanted.
  % Optional. Is for examiners receiving a revised version of the thesis to help them check that any required changes have been made. It is removed from the thesis before the final submitted version.
  % If this section is wanted, uncomment this section. (Keep TC:ignore and TC:endignore commented, otherwise these will contribute to the automatic word count)
  % Most people probably find this easiest to make their list of revisions as a seperate document, but you can do it here if you would like. This would let you do a \pageref{sec:my_label} to a \label{sec:my_label} so can say e.g. "On page 14 I added more info", and the 14 be automatically calculated.
  % Policy asks for this to be removed from the final thesis, which messes up the page numbers somewhat. (Which is partly why it's likely easier to do the list of revisions as a seperate document.) Page numbers are done here so that the final post-award thesis is correct. In the pre-examination thesis the numbers displayed on the page will be one lower than the number displayed by the PDF reader. Remove this page if not wanted/needed, and/or is a first or second year report


  %\lipsum[1-3] % generate dummy text for here
  
%\end{uomlotr} 
%TC:endignore % turn word count back on



%%%%%%%%%%%%%%%%%% TITLE PAGE %%%%%%%%%%%%%%%%%%

% Title and author are automatically taken from the document meta-data defined above
\title{A data reduction algorithm incorporating a low power continuous wavelet transform for use in wearable electroencephalography systems}
\author{Alexander J.\ Casson}

% Set the below yourself
\faculty{Science and Engineering}                  % "Faculty of" is added automatically
\school{School of Engineering} % "School of" not added automatically (as if in AMBS, School comes at the end rather than the start), so enter full name of School here
\departmentordivision{Department of Electrical and Electronic Engineering} % "Department of" or "Division of" not added automatically (as some Schools use Departments while others use Divisions), so enter full name of Dept./Division here
\submitdate{2026}                                  % regulations ask only for the year, not month
\wordcount{\mywordcount}		                   % use \wordcount{} to set the count. Can just type in a number (e.g. \wordcount{1000}	if don't want to use the automatic count.) This is automatically displayed after the table of contents. 
\maketitle

% (Optional) Add keywords to the PDF, if wanted
\hypersetup{pdfkeywords={Keyword 1, Keyword 2, Keyword 3}}



%%%%%%%%%%%%%%%%%% LISTS OF CONTENT %%%%%%%%%%%%%%%%%%

% Probably don't need all of these unless final thesis
\uomtoc % contents 
\uomlof % figures
\uomlot % tables
\begin{uomlop} % list of publications.
  % Can use biblatex to \printbibliography[heading=none] to populate automatically, or can add a custom list via the tocloft package, but probably easier to just type in unless you have lots!
  Publications go here.
\end{uomlop}
\begin{uomterms}
  Enter terms and abbreviations as table or similar.
  % Can do just by hand, making a list of common terms.

  % TAGGING: Check whether this still works
  % Alternatively can use the glossaries package. For this:
  % Uncomment and add the below to the preamble (anywhere before \begin{document}) with as many terms as you'd like:
    % \usepackage[acronym,nonumberlist,section=section,nogroupskip]{glossaries}
    % \makeglossaries
    % \newacronym{gcd}{GCD}{Greatest Common Divisor}
    % \newacronym{lcm}{LCM}{Least Common Multiple}
  % Then inside \begin{uomterms} add:
    % \printglossary[type=\acronymtype,title=]
  % Then inside your main text have at least one use of an acronym:
    % \gls{gcd}
  % You might need to select "Recompile from scratch" (from the down arrow next to "Recompile" in Overleaf) for this to take effect.
  % This still leaves a large-ish blank space before it displays the terms for some reason, but is probably good enough. I'll look to fix this at some point
\end{uomterms}



%%%%%%%%%%%%%%%%%% ABSTRACT %%%%%%%%%%%%%%%%%%
\begin{abstract} % put abstract here. Limit is 1 page.
  This is abstract text. 
  
  \lipsum[1-2]
\end{abstract}%
\clearpage



%%%%%%%%%%%%%%%%%% LAY ABSTRACT %%%%%%%%%%%%%%%%%%
\begin{uomlay} % Optional. Delete if not wanted
  Optional. If desired, this is lay abstract text. 
  
  \lipsum[1-2]
\end{uomlay}



%%%%%%%%%%%%%%%%%% DECLARATION OF ORIGINALITY %%%%%%%%%%%%%%%%%%
\begin{uomoriginality} % Don't need unless final thesis.
  \uomoriginalitydeclaration 
  % If the standard originality decalaration is sufficient, saying no portion of the work has been submitted in support of an application for another degree or qualification, the above command will automatically add the required text and nothing else is needed in this section.
  % If the standard statment isn't sufficient, then comment out the \uomoriginalitydeclaration command and type in your own text here explaining the authorship of any re-used portions. 
\end{uomoriginality}



%%%%%%%%%%%%%%%%%% COPYRIGHT STATEMENT %%%%%%%%%%%%%%%%%%
\uomcopyrightstatement % Needed for thesis. Not needed for first and second year reports. No options are needed. Having this command will add the required declarations



%%%%%%%%%%%%%%%%%% ACKNOWLEDGEMENTS %%%%%%%%%%%%%%%%%%
\begin{uomacknowledgements} % Optional. Delete if not wanted. Probably don't need unless final thesis
Optional. If desired, acknowledgements go here.
\end{uomacknowledgements}



%%%%%%%%%%%%%%%%%% ABOUT THE AUTHOR %%%%%%%%%%%%%%%%%%
\begin{uomauthor} % Optional. Delete if not wanted
Optional. If desired, a brief statement for External Examiners giving the candidate’s degree(s) and research experience, even if the latter consists only of the work done for this thesis.
\end{uomauthor} 



%%%%%%%%%%%%%%%%%% AI DECLARATION %%%%%%%%%%%%%%%%%%
\begin{uomai} % Optional. Delete if not wanted
Generative Artificial Intelligence (AI) disclosure: I used [AI tool name] to assist in idea generation, image creation, and for feedback on grammar and content. I implemented some of its recommendations. I used [AI tool name] to explore ideas for visuals (one of which is used and cited on page~\pageref{sec:content}). % change/add more as needed. 
\end{uomai}



%%%%%%%%%%%%%%%%%% CONTENT NOTIFICATION %%%%%%%%%%%%%%%%%%
\begin{uomcontent} % Optional. Delete if not wanted
Optional. If desired, see dedicated policy at \url{https://documents.manchester.ac.uk/display.aspx?DocID=74816} for whether a content or trigger warning is required. Examples of such warnings are given in that policy.
\end{uomcontent} 



%%%%%%%%%%%%%%%%%% Start content %%%%%%%%%%%%%%%%%%
\uomstartmainbody % Don't delete. used to flag to the hyperlinks in the PDF that the main content is  starting



%%%%%%%%%%%%%%%%%% CHAPTER 1 %%%%%%%%%%%%%%%%%%
\chapter{Introduction} % probably better to use \input{} or \include{} for each chapter to draw content from other files rather than having everything in one big file
  
  \lipsum[1-5] % generate dummy text for here



%%%%%%%%%%%%%%%%%% CHAPTER 2 %%%%%%%%%%%%%%%%%%
\chapter{Literature review}

  \section{Introduction}
    \lipsum[1] 
  
  \section{Example display items}
    This is an example of providing a cross-reference to \ref{sec:really_good_work}. Similarly, this is an example cross-reference to a sub-section, \ref{sec:content}. This is an example cross-reference to an appendix, \ref{sec:example_appendix}.
    
    This is an example of adding references \cite{ref:jCAS09,ref:jCAS09a,ref:jCAS10}. If you want the author name, or similar, you can use: \citeauthor*{ref:jCAS09} in \citeyear{ref:jCAS09} introduced a really good idea. (This is for when you primarily use numbered citations, but occasionally need an author's name. If using author names as the reference everywhere, change style=ieee in the biblatex setup above to whatever reference style you want, and then just use the cite command.)

    For adding \emph{emphasis} use the emph command. In general, try to avoid texttt, textit, textbf and similar commands --- these don't tell the document \emph{why} the style is being changed, and this information is needed for screen readers. In contrast, emph gives meaning to the change in format. By default it uses italics, but you can change the style if you want to. 
    
    Here are two example tables. \ref{table:example_tabularx} is an example using tabularx to fix the size of the table to be the same width as the page, with columns that auto-wrap if the text is too long. Note this is currently blank. tabularx is not currently supported for tagging. If you need this, you need to use the other (non-accessible) template.
    
    \ref{table:example_tabular} is an example using tabular where the width of the table is set by its contents. It may go beyond the width of the page.
    \begin{table}
      \centering
      \caption[Short caption for list of tables]{Example table. Full caption goes here. Often a short caption in [] is used as well as the main caption to keep the list of figures tidy; it gets messy if there are long captions going over more than one line.}
      \label{table:example_tabularx}
% TAGGING: Re-add when tabularx supports tagging
%      \begin{tabularx}{\linewidth}{ZZZZZ}
%        \toprule
%        \multirow{2}{*}{Participant} & \multicolumn{2}{c}{Number / \%} & \multicolumn{2}{c}{Duration / \%} \\
%                                     & Prime dresses & Non-prime dresses & Prime dresses & Non-prime dresses \\
%	  \otoprule
%        1  & 33.33 & 33.91 & 20.83 & 18.42 \\
%        2  & 13.04 & 17.50 & 04.93 & 07.62 \\
%        3  & 22.73 & 20.10 & 13.00 & 08.20 \\
%        4  & 31.34 & 21.88 & 10.57 & 11.09 \\
%        5  & 08.47 & 19.32 & 03.04 & 09.73 \\
%        \hline
%        Mean & 16.4 & 16.5 & 07.8 & 07.5 \\
%        Standard deviation & 09.7 & 06.6 & 05.4 & 03.3 \\
%        \bottomrule
%      \end{tabularx}
    \end{table}

    \begin{table}
      \centering
      \caption[Short caption for list of tables]{Example table. Full caption goes here. Often a short caption in [] is used as well as the main caption to keep the list of figures tidy; it gets messy if there are long captions going over more than one line.}
      \label{table:example_tabular}
      \tagpdfsetup{table/header-rows={2}}
      \begin{tabular}{ccccc}
        \toprule
        \multirow{2}{*}{Participant} & \multicolumn{2}{c}{Number / \%} & \multicolumn{2}{c}{Duration / \%} \\
                                     & Prime dresses & Non-prime dresses & Prime dresses & Non-prime dresses \\
	  \otoprule
        1  & 33.33 & 33.91 & 20.83 & 18.42 \\
        2  & 13.04 & 17.50 & 04.93 & 07.62 \\
        3  & 22.73 & 20.10 & 13.00 & 08.20 \\
        4  & 31.34 & 21.88 & 10.57 & 11.09 \\
        5  & 08.47 & 19.32 & 03.04 & 09.73 \\
        \hline
        Mean & 16.4 & 16.5 & 07.8 & 07.5 \\
        Standard deviation & 09.7 & 06.6 & 05.4 & 03.3 \\
        \bottomrule
      \end{tabular}
    \end{table}
  
    This is an example equation in text \(2\sin{\omega t}\). Below is an example of a displayed equation. See \ref{equ:example} for info.
    \begin{equation}
        a^{2} + b^{2} = c^{2} \label{equ:example}
    \end{equation}

    Note that numbers are displayed differently in the text depending on how they are entered. Compare for example 123456 vs.\ \(123456\) vs.\ \num{123456}. Entering numbers directly, such as 1955, should be used for \emph{text mode} numbers. That is, those representing text (dates, page numbers, and similar). Numbers representing maths, or variables or similar, should be entered inside \textbackslash( \textbackslash) 
    % TAGGING: Re-add when supported for tagging
    % or \lstinline{\num} % Re-add when supported for tagging
    so they are typeset in the same way as they appear in an equation. (This requires a bit of discipline, but helps ensure consistent use of number styles throughout.) 
    % TAGGING: Re-add when supported for tagging
    % \lstinline{\num} is provided by the \lstinline{siunitx} package. The display it provides can be customised if wanted. Documentation is at \url{https://texdoc.org/serve/siunitx/0}.
    You may see some \LaTeX\ information where they put maths in \$ \$. Don't do this, use \textbackslash( \textbackslash). \$ \$ will work at the moment, but won't in the future.

    This is an example of a quote in text \textcquote{ref:jCAS10}{The electroencephalogram (EEG) is a classic non-invasive method for measuring a person's brainwaves}. Below is an example of a displayed quote. Both are using the csquotes package.
    \begin{displaycquote}{ref:jCAS10}
      Electrodes are placed on the scalp to detect the microvolt-sized signals that result from synchronized neuronal activity within the brain.
    \end{displaycquote}

    This is an example figure. See \ref{fig:uom_logo} for more details. Alternatively, can put multiple figures into a subfigure environment, for example \ref{fig:uom_logo_in_subfig}, or \ref{fig:subfig_b} to cite a particular sub-figure.
    \begin{figure}
      \centering
      \includegraphics[alt={Put short description for screen readers here},width=0.3\textwidth,keepaspectratio=true]{uom_logo.pdf}
      \caption[Short caption for list of figures]{Example figure. Full caption goes here. Often a short caption in [] is used as well as the main caption to keep the list of figures tidy; it gets messy if there are long captions going over more than one line.}
      \label{fig:uom_logo}
    \end{figure} 

    \begin{figure}
      \centering
      \begin{subfigure}{0.3\linewidth}
        \includegraphics[alt={Put short description for screen readers here},width=\textwidth,keepaspectratio=true]{uom_logo.pdf}
        \caption{}
        \label{fig:subfig_a}
      \end{subfigure}
      \begin{subfigure}{0.3\linewidth}
        \includegraphics[alt={Put short description for screen readers here},width=\textwidth,keepaspectratio=true]{uom_logo.pdf}
        \caption{}
        \label{fig:subfig_b}
      \end{subfigure}
      \begin{subfigure}{0.3\linewidth}
        \includegraphics[alt={Put short description for screen readers here},width=\textwidth,keepaspectratio=true]{uom_logo.pdf}
        \caption{}
        \label{fig:subfig_c}
     \end{subfigure}
     \caption{Three copies of the University logo. (a) Copy one. (b) Copy two. (c) Copy three.}
     \label{fig:uom_logo_in_subfig}
    \end{figure}

    Code listing is not currently supported for tagging. If you need this, you need to use the other (non-accessible) template.
    % TAGGING: Re-add when supported for tagging
    %An example code listing is given in \ref{code:example}. See \url{https://www.overleaf.com/learn/latex/Code_listing} for more examples of how to display code. Code in the body of the text can be included as \lstinline{for} or \lstinline{while} or \lstinline{main}.

    % TAGGING: Re-add when supported for tagging
    %\begin{lstlisting}[language=Python, caption={My code example}, label={code:example}]
%import numpy as np
    
%def my_filter(in,f_obj):
%    y = filter(f_obj,in)
    
%    return y
%    \end{lstlisting}
  
  \section{Summary}
    \lipsum[6] % adds dummy text to fill up the example
  
  
  
%%%%%%%%%%%%%%%%%% CHAPTER 3 %%%%%%%%%%%%%%%%%%
\chapter{Really good work} \label{sec:really_good_work}

  \section{Introduction}
    \lipsum[1] % adds dummy text to fill up the example
  
  \section{Content} \label{sec:content}
    \subsection{Introduction}
	\lipsum[1] % adds dummy text to fill up the example
	
	\subsection{Detail}
        \lipsum[7-11] % adds dummy text to fill up the example
    
	\subsection{More detail}
	\lipsum[1-3] % adds dummy text to fill up the example
	
	\subsection{Summary}
	\lipsum[1] % adds dummy text to fill up the example
  
  \section{Summary}
    \lipsum[6] % adds dummy text to fill up the example



%%%%%%%%%%%%%%%%%% CONCLUSIONS %%%%%%%%%%%%%%%%%%
\chapter{Conclusions}
  \lipsum[1] % adds dummy text to fill up the example



%%%%%%%%%%%%%%%%%% REFERENCES %%%%%%%%%%%%%%%%%%
\printbibliography[title={References},heading=bibintoc] % a single list of references for the whole thesis



%%%%%%%%%%%%%%%%%% APPENDICES %%%%%%%%%%%%%%%%%%
\begin{uomappendix} 
  \chapter{First Appendix} \label{sec:example_appendix}
    \section{Section in Appendix}
    \lipsum[1-6]

  \chapter{Second Appendix}
    \section{Section in Appendix}
    \lipsum[1-6]
\end{uomappendix}



%%%%%%%%%%%%%%%%%% END MATTER %%%%%%%%%%%%%%%%%%
\end{document}